\documentclass[11pt,a4paper]{article}

\usepackage[a4paper,margin=2.45cm]{geometry}
\usepackage[T1]{fontenc}
\usepackage[utf8]{inputenc}
\usepackage{lmodern}
\usepackage{amsmath,amssymb,amsthm,mathtools}
\usepackage{enumitem}
\usepackage{booktabs}
\usepackage{array}
\usepackage{xcolor}
\usepackage{microtype}
\usepackage{fancyhdr}
\usepackage{hyperref}
\usepackage{bookmark}

\hypersetup{
  colorlinks=true,
  linkcolor=blue!55!black,
  urlcolor=blue!55!black,
  citecolor=blue!55!black,
  pdftitle={Simple Homotopy Theory and Manifold Topology: Revised Lecture Notes},
  pdfauthor={OpenAI}
}

\pagestyle{fancy}
\fancyhf{}
\lhead{Simple Homotopy Theory and Manifold Topology}
\rhead{Revised Notes}
\cfoot{\thepage}
\setlength{\headheight}{14pt}

\setlist[itemize]{leftmargin=2em,itemsep=0.25em,topsep=0.35em}
\setlist[enumerate]{leftmargin=2.2em,itemsep=0.25em,topsep=0.35em}
\linespread{1.08}

\newcommand{\Wh}{\operatorname{Wh}}
\newcommand{\GL}{\operatorname{GL}}
\newcommand{\im}{\operatorname{im}}
\newcommand{\Cone}{\operatorname{Cone}}
\newcommand{\rel}{\operatorname{rel}}
\newcommand{\ab}{\operatorname{ab}}
\newcommand{\ZZ}{\mathbb{Z}}
\newcommand{\id}{\operatorname{id}}
\newcommand{\widet}{\widetilde}

\theoremstyle{definition}
\newtheorem{definition}{Definition}[section]
\newtheorem{example}[definition]{Example}
\newtheorem{remark}[definition]{Remark}
\newtheorem{convention}[definition]{Convention}

\theoremstyle{plain}
\newtheorem{theorem}[definition]{Theorem}
\newtheorem{proposition}[definition]{Proposition}
\newtheorem{corollary}[definition]{Corollary}

\theoremstyle{remark}
\newtheorem*{proofidea}{Idea of proof}
\newtheorem*{note}{Note}

\title{Simple Homotopy Theory and Manifold Topology\\\large Revised Lecture Notes}
\author{}
\date{June 16, 2026}

\begin{document}
\maketitle

\begin{center}
\begin{minipage}{0.88\textwidth}
These notes review simple homotopy theory, Whitehead torsion, and the smooth \(s\)-cobordism theorem. The exposition follows the original lecture order: elementary moves on finite CW complexes, the appearance of group rings and \(K_1\), the definition and properties of Whitehead torsion, and finally the role of torsion in high-dimensional manifold topology.
\end{minipage}
\end{center}

\tableofcontents
\newpage

\part*{Part I. Revised Lecture Notes}
\addcontentsline{toc}{part}{Part I. Revised Lecture Notes}

\section{Finite CW complexes and elementary moves}

We begin with the cellular language used throughout the notes. A CW complex is built inductively by attaching cells. The \(n\)-skeleton \(X^n\) is obtained from \(X^{n-1}\) by attaching \(n\)-disks along maps
\[
  \varphi_\alpha:S^{n-1}\longrightarrow X^{n-1}.
\]
Equivalently,
\[
  X^n=X^{n-1}\cup_{\coprod_\alpha \varphi_\alpha}\coprod_\alpha D^n_\alpha .
\]
A CW complex is called \emph{finite} if it has only finitely many cells. Finite CW complexes are the natural setting for simple homotopy theory because the elementary operations below are finite combinatorial operations.

\begin{convention}[Regularity for free faces]
The language of a ``free face'' is most transparent for finite regular CW complexes or finite cellular complexes with a combinatorial structure. In a completely general CW complex, the word ``face'' requires extra care because characteristic maps need not be embeddings on closed cells. In these notes, elementary collapses are understood in a setting where deleting the two cells in question leaves a genuine subcomplex.
\end{convention}

\subsection{Elementary collapses}

\begin{definition}[Free face]
Let \(K\) be a finite regular CW complex. Suppose \(e^{n-1}\) is an \((n-1)\)-cell and \(e^n\) is an \(n\)-cell of \(K\). We say that \(e^{n-1}\) is a \emph{free face} of \(e^n\) if:
\begin{enumerate}
  \item \(e^{n-1}\) occurs in the boundary of \(e^n\) in the standard cellular way; and
  \item \(e^{n-1}\) is not contained in the boundary of any other \(n\)-cell, so that removing \(e^{n-1}\) together with \(e^n\) leaves a subcomplex.
\end{enumerate}
The pair \((e^{n-1},e^n)\) is called a \emph{free pair}.
\end{definition}

Intuitively, the pair \((e^{n-1},e^n)\) is redundant. The higher-dimensional cell \(e^n\) is attached along a boundary containing a face exposed to the outside, so the two cells can be removed without changing the ordinary homotopy type.

\begin{definition}[Elementary collapse]
Let \(K\) be a finite regular CW complex and let \((e^{n-1},e^n)\) be a free pair. Let \(L\) be the subcomplex obtained by deleting the interiors of these two cells. We say that \(K\) \emph{elementarily collapses} to \(L\), and write
\[
  K\searrow L.
\]
A finite composition of elementary collapses is called a \emph{collapse}.
\end{definition}

\begin{proposition}
If \(K\searrow L\) is an elementary collapse, then the inclusion
\[
  L\hookrightarrow K
\]
is a homotopy equivalence. In fact, \(K\) deformation retracts onto \(L\).
\end{proposition}

\begin{proofidea}
The free pair is locally modeled on a disk attached along all but one exposed face. Since the exposed face is not used by any other cell, one can push points of \(e^n\cup e^{n-1}\) away from the free face and into the remaining part of the boundary. This local deformation is the identity on \(L\), so it extends to a deformation retraction \(K\to L\).
\end{proofidea}

\subsection{Elementary expansions}

\begin{definition}[Elementary expansion]
An \emph{elementary expansion} is the inverse operation to an elementary collapse. If \(K\searrow L\), then \(K\) is obtained from \(L\) by adding a pair of cells \((e^{n-1},e^n)\) in such a way that \(e^{n-1}\) becomes a free face of \(e^n\) after the addition.
\end{definition}

Both elementary collapses and elementary expansions preserve ordinary homotopy type. Simple homotopy theory is finer than ordinary homotopy theory because it records whether a homotopy equivalence can be generated by these elementary cellular moves.

\subsection{Simple homotopy equivalence}

\begin{definition}[Simple homotopy equivalence]
Two finite CW complexes \(X\) and \(Y\) are \emph{simple homotopy equivalent}, written
\[
  X\simeq_s Y,
\]
if there exists a finite zig-zag
\[
  X=K_0\longleftrightarrow K_1\longleftrightarrow\cdots\longleftrightarrow K_r=Y,
\]
where each arrow is an elementary expansion or an elementary collapse.

A homotopy equivalence \(f:X\to Y\) is called a \emph{simple homotopy equivalence} if, up to homotopy, it is represented by the homotopy equivalence associated with such a finite sequence of elementary expansions and collapses.
\end{definition}

\begin{remark}
The word ``simple'' does not mean that the spaces themselves are geometrically simple. It means that the equivalence between them is generated by elementary cellular operations.
\end{remark}

\begin{example}[Contractible versus collapsible]
A finite CW complex \(K\) is \emph{contractible} if \(K\simeq *\). It is \emph{collapsible} if there is a sequence of elementary collapses
\[
  K\searrow \cdots\searrow *.
\]
Every collapsible complex is contractible, but not every contractible finite CW complex is collapsible. This is one of the simplest indications that simple homotopy theory is more refined than ordinary homotopy theory.
\end{example}

\section{The geometric Whitehead group}

The previous section suggests a natural question: given a homotopy equivalence
\[
  f:X\longrightarrow Y
\]
of finite CW complexes, how can one decide whether it is simple? Whitehead's answer is that there is an obstruction group and a torsion element in that group.

One can first describe this obstruction geometrically. Fix a finite connected CW complex \(X\). Consider pairs \((K,X)\), where \(K\) is a finite CW complex containing \(X\) as a subcomplex and the inclusion
\[
  X\hookrightarrow K
\]
is a homotopy equivalence. Such a pair can be regarded as a finite extension of \(X\) which does not change its ordinary homotopy type.

If \(K\) is obtained from \(X\) by elementary expansions, then the inclusion is visibly simple. In general, however, \(K\) may be homotopy equivalent to \(X\) without being reducible to \(X\) by an obvious sequence of elementary collapses. The geometric Whitehead group records this remaining obstruction.

The full construction requires a stabilization relation and a composition law. For these notes, the essential statement is the following.

\begin{theorem}[Geometric form of Whitehead torsion]
There is an abelian group, denoted geometrically by \(\Wh(X)\), such that every homotopy equivalence of finite connected CW complexes
\[
  f:X\longrightarrow Y
\]
has an associated torsion element
\[
  \tau^{\operatorname{top}}(f)\in \Wh(Y).
\]
Moreover,
\[
  \tau^{\operatorname{top}}(f)=0
  \quad\Longleftrightarrow\quad
  f \text{ is a simple homotopy equivalence}.
\]
\end{theorem}

This statement already gives the main philosophy. To see why the fundamental group appears, however, one must translate the geometry into algebra.

\section{Algebraic input: group rings and \texorpdfstring{\(K_1\)}{K1}}

\subsection{The group ring}

Let \(G\) be a group. The integral group ring \(\ZZ[G]\) consists of finite formal sums
\[
  \sum_{g\in G} n_g g,\qquad n_g\in\ZZ,
\]
with addition defined coefficientwise and multiplication determined by the group operation in \(G\), extended linearly. If \(G\) is nonabelian, then \(\ZZ[G]\) is usually noncommutative.

The group ring enters topology through universal covers. If \(X\) is connected and \(\widetilde X\to X\) is its universal cover, then \(\pi_1X\) acts on \(\widetilde X\) by deck transformations. Hence the cellular chain groups
\[
  C_n(\widetilde X)
\]
are naturally modules over \(\ZZ[\pi_1X]\). Once one chooses a lift of each cell of \(X\), the module \(C_n(\widetilde X)\) becomes a finitely generated free based \(\ZZ[\pi_1X]\)-module.

Thus a finite CW complex gives a finite based free chain complex over the group ring:
\[
  \cdots\longrightarrow C_n(\widetilde X)
  \xrightarrow{\partial_n} C_{n-1}(\widetilde X)
  \longrightarrow\cdots\longrightarrow C_0(\widetilde X).
\]
The boundary maps are represented, after choices of bases, by matrices over \(\ZZ[\pi_1X]\).

\begin{remark}
Some authors use left \(\ZZ[\pi_1X]\)-modules and others use right modules. This changes matrix conventions but not the underlying torsion invariant once a convention is fixed.
\end{remark}

\subsection{The group \texorpdfstring{\(K_1(R)\)}{K1(R)}}

Let \(R\) be a unital associative ring. The group \(K_1(R)\) is constructed from invertible matrices over \(R\). Let
\[
  \GL_n(R)
\]
denote the group of invertible \(n\times n\) matrices over \(R\). The stabilization maps
\[
  \GL_n(R)\longrightarrow \GL_{n+1}(R),
  \qquad
  A\longmapsto
  \begin{pmatrix}
    A&0\\
    0&1
  \end{pmatrix}
\]
allow one to form the stable general linear group
\[
  \GL(R)=\bigcup_{n\ge 1}\GL_n(R).
\]

\begin{definition}[Algebraic \(K_1\)]
The group \(K_1(R)\) is the abelianization of \(\GL(R)\):
\[
  K_1(R)=\GL(R)_{\ab}.
\]
Equivalently, one quotients \(\GL(R)\) by its commutator subgroup.
\end{definition}

When \(R\) is commutative, the determinant gives a familiar homomorphism from \(K_1(R)\) to \(R^\times\). For noncommutative rings, the ordinary determinant is not available in the same way, and \(K_1(R)\) functions as the correct determinant-like invariant.

\subsection{The Whitehead group}

For \(R=\ZZ[G]\), there are obvious invertible elements: for each \(g\in G\), both \(g\) and \(-g\) are units in \(\ZZ[G]\). These are called the \emph{trivial units}. They arise topologically from choices of lifts and orientations of cells, so they should not contribute to a topological invariant.

\begin{definition}[Whitehead group]
Let \(G\) be a group. The Whitehead group of \(G\) is
\[
  \Wh(G)=K_1(\ZZ[G])\big/\langle [\pm g]\mid g\in G\rangle,
\]
where \([\pm g]\) denotes the class in \(K_1(\ZZ[G])\) of the \(1\times 1\) matrix whose entry is \(\pm g\). If \(X\) is a connected CW complex, we write
\[
  \Wh(X)=\Wh(\pi_1X).
\]
\end{definition}

\begin{remark}
The quotient by the classes \([\pm g]\) is forced by topology. Choosing a different lift of a cell in the universal cover multiplies a basis vector by an element of \(\pi_1X\). Reversing the orientation of a cell multiplies a basis vector by \(-1\). Whitehead torsion must be independent of these choices, so these units are divided out.
\end{remark}

\begin{example}
If \(G=1\), then \(\ZZ[G]=\ZZ\). Since \(K_1(\ZZ)\cong\{\pm 1\}\), quotienting by the trivial units gives
\[
  \Wh(1)=0.
\]
This fact is crucial in the simply connected case of the \(h\)-cobordism theorem.
\end{example}

\section{Whitehead torsion of chain equivalences}

We now explain how torsion is extracted from chain complexes. The goal is not to reproduce every technical detail, but to make clear why a homotopy equivalence of CW complexes has an associated class in \(\Wh(\pi_1)\).

\subsection{Based free chain complexes}

Let \(R\) be a ring. A finite based free chain complex over \(R\) is a finite chain complex
\[
  0\longrightarrow C_n\xrightarrow{\partial_n} C_{n-1}
  \longrightarrow\cdots\longrightarrow C_0\longrightarrow 0,
\]
where each \(C_i\) is a finitely generated free \(R\)-module equipped with a chosen basis.

If such a complex is acyclic and the relevant exact sequences are split in the way needed to choose bases, one can compare two kinds of bases: the original cellular bases of the chain groups and the bases assembled from boundaries and chosen lifts of boundaries. The resulting change-of-basis matrices determine an element of \(K_1(R)\). This is the torsion of the complex.

\begin{definition}[Torsion of an acyclic based complex, schematic form]
Let \(C_*\) be a finite based free acyclic chain complex over \(R\). Choose bases for the boundary modules \(B_i=\im(\partial_{i+1})\), together with lifts of these bases into the preceding chain groups. In each degree, these choices produce a new basis of \(C_i\). Comparing this new basis with the original basis gives an invertible matrix over \(R\). The alternating product of their classes defines the torsion
\[
  \tau(C_*)\in K_1(R),
\]
up to the standard choices and sign conventions.
\end{definition}

\begin{remark}[A technical point]
Over a general noncommutative ring, the statement ``choose bases for all boundary modules'' should be treated with care: submodules of free modules need not be free. In the standard definition of Whitehead torsion, one works with finite based free complexes that are contractible, or uses equivalent formulations involving a chain contraction. For cellular chain complexes over \(\ZZ[\pi]\), the construction is arranged so that the final class in the Whitehead group is independent of the auxiliary choices relevant to topology.
\end{remark}

The exact sign convention varies between authors. What matters here is the principle: a suitably based contractible or acyclic chain complex has a determinant-like invariant, and for cellular complexes over \(\ZZ[\pi_1]\) this invariant naturally lands in a Whitehead group after quotienting by the trivial units.

\subsection{The mapping cone construction}

Let
\[
  \phi:C_*\longrightarrow D_*
\]
be a chain map between finite based free chain complexes over \(R\). Its mapping cone is the chain complex \(\Cone(\phi)_*\) defined by
\[
  \Cone(\phi)_n=D_n\oplus C_{n-1},
\]
with differential
\[
  \partial_{\Cone}(d,c)=\bigl(\partial_D d+\phi(c),-\partial_C c\bigr).
\]

A standard theorem says that \(\phi\) is a chain homotopy equivalence if and only if \(\Cone(\phi)_*\) is contractible. Acyclicity alone is weaker over a general ring: an acyclic cone means that \(\phi\) is a quasi-isomorphism, but it need not imply chain homotopy equivalence without extra hypotheses.

Therefore, if \(\phi\) is a chain homotopy equivalence, one defines its torsion by taking the torsion of its mapping cone:
\[
  \tau(\phi):=\tau\bigl(\Cone(\phi)_*\bigr).
\]

\subsection{From a cellular map to torsion}

Let
\[
  f:X\longrightarrow Y
\]
be a cellular homotopy equivalence of finite connected CW complexes. After identifying the fundamental groups via \(f_*:\pi_1X\to\pi_1Y\), and after choosing a lift of \(f\) to the universal covers, \(f\) induces a chain homotopy equivalence
\[
  C_*(\widetilde X)\longrightarrow C_*(\widetilde Y)
\]
of finite based free \(\ZZ[\pi_1Y]\)-chain complexes. The torsion of this chain equivalence gives
\[
  \tau(f)\in \Wh(\pi_1Y).
\]

The use of \(\Wh(\pi_1Y)\), rather than the larger group \(K_1(\ZZ[\pi_1Y])\), reflects the fact that the construction should not depend on arbitrary choices of lifts and orientations of cells.

\section{The fundamental theorem of simple homotopy theory}

We can now state the central theorem connecting elementary cellular geometry with the algebraic torsion invariant.

\begin{theorem}[Whitehead]
Let \(f:X\to Y\) be a homotopy equivalence of finite connected CW complexes. Then
\[
  f \text{ is a simple homotopy equivalence}
  \quad\Longleftrightarrow\quad
  \tau(f)=0\in \Wh(\pi_1Y).
\]
\end{theorem}

\begin{proofidea}
There are two directions.

First, suppose \(f\) is simple. Then \(f\) is built from elementary expansions and collapses. On the chain level, an elementary expansion adds an elementary contractible summand to the cellular chain complex. Such a summand has trivial torsion. Since torsion is additive under composition, the total torsion of a finite simple sequence is zero.

Conversely, suppose \(\tau(f)=0\). One studies the mapping cylinder of \(f\) and the relative cellular chain complex associated to the relevant pair. By performing elementary cellular modifications, the relative boundary maps can be changed by stabilization and by elementary row and column operations. These operations match the algebraic operations used in defining \(K_1(\ZZ[\pi_1])\). Vanishing torsion means that the remaining matrix obstruction can be reduced to the identity by these operations. Geometrically, this reduction realizes \(f\) by elementary expansions and collapses. Hence \(f\) is simple.
\end{proofidea}

This theorem is the bridge between the combinatorial and algebraic sides of the theory. The geometric problem of recognizing simple homotopy equivalences is converted into an algebraic problem in \(K_1\) of a group ring.

\subsection{Basic properties of torsion}

Whitehead torsion satisfies structural properties that explain why it behaves like a determinant.

\begin{proposition}[Homotopy invariance]
If \(f,g:X\to Y\) are homotopic homotopy equivalences of finite CW complexes, then
\[
  \tau(f)=\tau(g).
\]
\end{proposition}

\begin{proposition}[Composition formula]
Let
\[
  X\xrightarrow{f}Y\xrightarrow{g}Z
\]
be homotopy equivalences of finite connected CW complexes. Then
\[
  \tau(g\circ f)=\tau(g)+g_*\tau(f),
\]
where \(g_*:\Wh(\pi_1Y)\to\Wh(\pi_1Z)\) is the isomorphism induced by \(g\).
\end{proposition}

\begin{corollary}
If \(f:X\to Y\) is a homotopy equivalence, then the torsion of a homotopy inverse is determined by \(\tau(f)\). More precisely, for a homotopy inverse \(h:Y\to X\),
\[
  \tau(h)=-h_*\tau(f).
\]
In particular, if \(\tau(f)=0\), then any homotopy inverse of \(f\) is also simple.
\end{corollary}

\section{Manifold topology and handle decompositions}

We now pass from CW complexes to manifolds. Handles are the smooth analogues of cells.

Let \(W\) be a compact smooth \((n+1)\)-manifold with boundary, and suppose \(M_0\subset\partial W\) is a boundary component. A handle decomposition of \(W\) relative to \(M_0\) describes \(W\) as being built from \(M_0\times[0,1]\) by attaching handles.

A \(k\)-handle in dimension \(n+1\) is a copy of
\[
  D^k\times D^{n+1-k}.
\]
It is attached along
\[
  S^{k-1}\times D^{n+1-k}.
\]
The integer \(k\) is the \emph{index} of the handle.

Morse theory gives a useful way to visualize this construction. A Morse function on \(W\) has critical points, and each critical point of index \(k\) contributes a \(k\)-handle. Simplifying a handle decomposition corresponds to simplifying the Morse function.

\subsection{Handle cancellation}

The smooth counterpart of an elementary collapse is handle cancellation. Roughly, a \(k\)-handle and a \((k+1)\)-handle cancel if the attaching sphere of the \((k+1)\)-handle intersects the belt sphere of the \(k\)-handle transversely in exactly one point, with the correct sign.

When such a cancellation is possible, the two handles can be removed from the decomposition without changing the diffeomorphism type of the manifold relative to the lower boundary.

\subsection{Handle slides and matrices}

In a handle decomposition, handles of the same index may be slid over one another. Algebraically, handle slides correspond to elementary row and column operations on matrices over the group ring \(\ZZ[\pi_1]\). Stabilizing a handle decomposition by adding a canceling pair corresponds to stabilizing a matrix by adding an identity block.

This is the geometric origin of the appearance of \(K_1(\ZZ[\pi_1])\) in the \(s\)-cobordism theorem. The algebraic operations that define Whitehead torsion are reflected precisely in the geometric operations available for handles.

\section{Cobordisms and \texorpdfstring{\(h\)}{h}-cobordisms}

\begin{definition}[Cobordism]
Let \(M_0\) and \(M_1\) be closed smooth \(n\)-manifolds. A \emph{cobordism} from \(M_0\) to \(M_1\) is a compact smooth \((n+1)\)-manifold \(W\) whose boundary is identified with the disjoint union
\[
  \partial W=M_0\sqcup M_1.
\]
We write this as \((W;M_0,M_1)\).
\end{definition}

\begin{definition}[\(h\)-cobordism]
A cobordism \((W;M_0,M_1)\) is an \emph{\(h\)-cobordism} if the inclusions
\[
  M_0\hookrightarrow W,
  \qquad
  M_1\hookrightarrow W
\]
are homotopy equivalences.
\end{definition}

The product cobordism
\[
  M_0\times[0,1]
\]
is the trivial example of an \(h\)-cobordism from \(M_0\) to itself. The main geometric question is whether every \(h\)-cobordism is diffeomorphic to such a product.

\begin{definition}[Torsion of an \(h\)-cobordism]
Let \((W;M_0,M_1)\) be an \(h\)-cobordism and let
\[
  i:M_0\hookrightarrow W
\]
be the inclusion. Its Whitehead torsion is first naturally an element of \(\Wh(\pi_1W)\). Since \(i\) is a homotopy equivalence, \(i_*:\pi_1M_0\to\pi_1W\) is an isomorphism, so we identify this group with \(\Wh(\pi_1M_0)\) and write
\[
  \tau(W,M_0):=\tau(i)\in \Wh(\pi_1M_0).
\]
\end{definition}

If \(W\cong M_0\times[0,1]\) relative to \(M_0\), then the inclusion \(M_0\hookrightarrow W\) is a simple homotopy equivalence. Hence
\[
  \tau(W,M_0)=0.
\]
The content of the \(s\)-cobordism theorem is the converse in high dimensions.

\section{The \texorpdfstring{\(s\)}{s}-cobordism theorem}

\begin{theorem}[Smooth \(s\)-cobordism theorem]
Let \((W;M_0,M_1)\) be a compact smooth \(h\)-cobordism between closed smooth manifolds. Assume
\[
  \dim W\ge 6.
\]
Then \(W\) is diffeomorphic to the product \(M_0\times[0,1]\) relative to \(M_0\) if and only if
\[
  \tau(W,M_0)=0\in \Wh(\pi_1M_0).
\]
Equivalently,
\[
  W\cong M_0\times[0,1]\ \rel\ M_0
  \quad\Longleftrightarrow\quad
  \tau(W,M_0)=0.
\]
\end{theorem}

The letter \(s\) stands for ``simple.'' An \(h\)-cobordism whose boundary inclusion is a simple homotopy equivalence is called an \emph{\(s\)-cobordism}. Thus, in dimensions at least six, simple \(h\)-cobordisms are precisely product cobordisms.

\subsection{Relation with the \texorpdfstring{\(h\)}{h}-cobordism theorem}

If \(M_0\) is simply connected, then
\[
  \Wh(\pi_1M_0)=\Wh(1)=0.
\]
Therefore every \(h\)-cobordism with simply connected boundary has zero torsion. The \(s\)-cobordism theorem then reduces to the classical \(h\)-cobordism theorem:
\[
  \text{simply connected high-dimensional \(h\)-cobordism}
  \quad\Longrightarrow\quad
  \text{product cobordism}.
\]
This is one reason the simply connected case is more accessible: there is no Whitehead torsion obstruction.

\subsection{Why the dimension assumption appears}

The proof uses transversality and the Whitney trick. The Whitney trick removes pairs of intersection points between submanifolds by finding embedded Whitney disks. In high dimensions there is enough room to arrange these disks and keep the relevant submanifolds disjoint during the required isotopies.

In low dimensions the same geometric moves can fail. Dimension four is especially subtle: embedded disks and intersection phenomena behave in ways that are not controlled by the same elementary high-dimensional arguments. Thus the hypothesis \(\dim W\ge 6\) is not cosmetic; it is built into the geometric proof.

\section{Proof strategy of the \texorpdfstring{\(s\)}{s}-cobordism theorem}

We outline the proof because it explains the role of Whitehead torsion most clearly.

\subsection{Step 1: Choose a handle decomposition}

Start with a handle decomposition of \(W\) relative to \(M_0\). Since \((W;M_0,M_1)\) is an \(h\)-cobordism, the relative homology groups
\[
  H_*(W,M_0)
\]
vanish. More strongly, after passing to the universal cover, the relative cellular or handle chain complex over \(\ZZ[\pi_1M_0]\) is contractible in the relevant simple-homotopy sense.

\subsection{Step 2: Move handles into middle dimensions}

Using handle cancellation below and above the middle dimension, together with transversality and the Whitney trick, one simplifies the handle decomposition. A normal form result says that, after suitable modifications, all remaining handles may be arranged to occur in two consecutive middle dimensions.

Schematically, the remaining relative handle chain complex has the form
\[
  0\longrightarrow C_{k+1}\xrightarrow{d}C_k\longrightarrow 0.
\]
The differential \(d\) is represented by a square matrix over \(\ZZ[\pi_1M_0]\).

\subsection{Step 3: Interpret the matrix}

Because the relative complex is contractible, the matrix representing \(d\) is invertible over the group ring after the appropriate choices. Its class in \(K_1(\ZZ[\pi_1M_0])\), modulo the trivial units \([\pm g]\), is exactly the Whitehead torsion
\[
  \tau(W,M_0)\in \Wh(\pi_1M_0).
\]

\subsection{Step 4: Use vanishing torsion}

If \(\tau(W,M_0)=0\), then the matrix can be reduced to the identity using stabilization and elementary row and column operations. Geometrically, these operations correspond to adding canceling handle pairs and performing handle slides.

After these modifications, the middle-dimensional handles are arranged in canceling pairs. Handle cancellation removes them. Therefore all handles disappear, and one obtains
\[
  W\cong M_0\times[0,1]\ \rel\ M_0.
\]

\subsection{Step 5: Prove the converse}

If \(W\cong M_0\times[0,1]\) relative to \(M_0\), then the inclusion \(M_0\hookrightarrow W\) is a simple homotopy equivalence. Hence its torsion is zero by the fundamental theorem of simple homotopy theory. This proves the equivalence in the theorem.

\section{Consequences and examples}

\subsection{The simply connected case}

As noted above, if \(\pi_1M_0=1\), then
\[
  \Wh(\pi_1M_0)=0.
\]
Thus every high-dimensional simply connected \(h\)-cobordism is a product. This is the classical \(h\)-cobordism theorem.

This theorem is one of the central tools behind the high-dimensional Poincar\'e conjecture. It shows that, in high dimensions, homotopy-theoretic information combined with suitable cobordism input can force a manifold to be standard in a strong geometric sense. In the smooth category, one must still distinguish homeomorphism from diffeomorphism, since exotic smooth structures may occur.

\subsection{The non-simply connected case}

When \(\pi_1M_0\) is nontrivial, the Whitehead group may be nonzero. In that case an \(h\)-cobordism can have nonzero torsion and therefore need not be a product. The \(s\)-cobordism theorem does not say that all \(h\)-cobordisms are trivial; it says exactly which ones are trivial.

In this sense, Whitehead torsion is the missing invariant in the non-simply connected \(h\)-cobordism problem.

\subsection{Classification viewpoint}

A refined form of the \(s\)-cobordism theorem says that, for a fixed closed smooth manifold \(M\) of dimension at least five, equivalence classes of smooth \(h\)-cobordisms starting from \(M\), up to diffeomorphism relative to \(M\), are classified by
\[
  \Wh(\pi_1M).
\]
Under this viewpoint, the product cobordism corresponds to the zero element. Nonzero elements of the Whitehead group correspond to nontrivial \(h\)-cobordisms.

\newpage
\part*{Part II. Revision Summary and Mathematical Notes}
\addcontentsline{toc}{part}{Part II. Revision Summary and Mathematical Notes}

\section*{Main revisions}
\addcontentsline{toc}{section}{Main revisions}

The revised version preserves the original mathematical order: elementary collapses and expansions, simple homotopy equivalence, group rings and \(K_1\), Whitehead torsion, handle decompositions, and the \(s\)-cobordism theorem. The main changes are as follows.

\begin{enumerate}
  \item The exposition has been reorganized into definitions, propositions, theorems, proof ideas, examples, and remarks, so that the notes are easier to use for review.
  \item The notation has been standardized. In particular, the universal cover is written as \(\widetilde X\), the Whitehead group as \(\Wh(G)\), the mapping cone as \(\Cone(\phi)\), and the relevant chain complexes are consistently treated as based free modules over group rings.
  \item Several explanatory transitions have been added, especially at the points where group rings enter through universal covers, where trivial units are quotiented out, and where handle slides correspond to elementary matrix operations.
  \item The proof sketches were made more explicit without expanding them into full proofs. The central theme is now clearer: elementary geometric operations correspond to algebraic matrix operations, and Whitehead torsion measures the remaining obstruction.
  \item The statement of the \(s\)-cobordism theorem was adjusted to emphasize the high-dimensional smooth setting and the meaning of the condition \(\dim W\ge 6\).
\end{enumerate}

\section*{Mathematical issues checked and corrected}
\addcontentsline{toc}{section}{Mathematical issues checked and corrected}

The main mathematical direction of the original notes was correct. A few points needed more precise formulation.

\begin{enumerate}
  \item \textbf{Free faces in CW complexes.} The term ``free face'' is completely natural for regular CW complexes or combinatorial cell complexes. In arbitrary CW complexes, one must be more careful because closed cells need not be embedded disks. The revised notes state the regularity convention explicitly.

  \item \textbf{Mapping cone criterion.} The original phrasing suggested that a chain map is a chain homotopy equivalence if and only if its mapping cone is acyclic. Over a general ring, the correct criterion is that the mapping cone is \emph{contractible}. Acyclicity alone is weaker and generally only detects quasi-isomorphisms. This has been corrected.

  \item \textbf{Definition of torsion for based complexes.} The intuitive definition using bases of boundary modules is useful pedagogically, but over a general group ring submodules of free modules need not be free. The revised notes keep the intuition while adding the needed technical warning and connecting the construction to contractible based free complexes or chain contractions.

  \item \textbf{Whitehead group quotient.} The quotient in the definition of \(\Wh(G)\) is more precisely taken by the subgroup of \(K_1(\ZZ[G])\) generated by the classes of the \(1\times1\) matrices \([\pm g]\), where \(g\in G\). The revised notes write
  \[
    \Wh(G)=K_1(\ZZ[G])\big/\langle [\pm g]\mid g\in G\rangle.
  \]

  \item \textbf{Target of the torsion of an \(h\)-cobordism.} For the inclusion \(i:M_0\hookrightarrow W\), the torsion naturally lies in \(\Wh(\pi_1W)\). Since \(i\) is a homotopy equivalence, one identifies \(\pi_1M_0\cong\pi_1W\), and hence regards the torsion as an element of \(\Wh(\pi_1M_0)\). This identification is now stated explicitly.

  \item \textbf{Classification viewpoint.} The statement that \(h\)-cobordisms are classified by \(\Wh(\pi_1M)\) requires high-dimensional hypotheses and a fixed category, such as the smooth category used here. The revised notes specify the smooth high-dimensional setting and the equivalence relation relative to the initial boundary.
\end{enumerate}

\section*{References}
\addcontentsline{toc}{section}{References}

\begin{enumerate}[label={[\arabic*]}]
  \item M. M. Cohen, \emph{A Course in Simple-Homotopy Theory}, Graduate Texts in Mathematics, Vol. 10, Springer, 1973.
  \item J. F. Davis and P. Kirk, \emph{Lecture Notes in Algebraic Topology}, Graduate Studies in Mathematics, Vol. 35, American Mathematical Society, 2001.
  \item A. Hatcher, \emph{Algebraic Topology}, Cambridge University Press, 2002.
  \item W. L\"uck and T. Macko, \emph{Surgery Theory: Foundations}, Grundlehren der mathematischen Wissenschaften, Vol. 362, Springer, 2024.
  \item J. Milnor, \emph{Lectures on the \(h\)-Cobordism Theorem}, Princeton University Press, 1965.
  \item F. Waldhausen, \emph{Algebraische Topologie II}, lecture notes.
\end{enumerate}

\end{document}
